%HOMEWORK template for M 325K
\documentclass{article}
\headheight=7pt         \topmargin=14pt
\textheight=574pt       \textwidth=445pt
\oddsidemargin=18pt     \evensidemargin=18pt

%Begin package import

\usepackage{amsmath, amssymb, amsthm}
\usepackage{mathtools}
\usepackage{pdfpages}
\usepackage{tikz-cd}

%End package import
%Begin custom commands

\newcommand*{\T}{\mathcal{T}}
\newcommand*{\HC}{\mathcal{H}}
\newcommand*{\B}{\mathcal{B}}
\newcommand*{\U}{\mathcal{U}}
\newcommand*{\cl}{\mathrm{cl}}
\newcommand*{\subeq}{\subseteq}
\newcommand*{\empset}{\emptyset}
\newcommand{\C}{\mathbb{C}}
\newcommand{\R}{\mathbb{R}}
\newcommand{\Q}{\mathbb{Q}}
\newcommand{\Z}{\mathbb{Z}}
\newcommand{\N}{\mathbb{N}}
\newcommand{\F}{\mathbb{F}}
\newcommand{\abs}[1]{\left\lvert #1\right\rvert}
\newcommand{\RM}[1]{\mathrm{#1}}
\newcommand{\BF}[1]{\textbf{#1}}
\newcommand{\e}{\epsilon}
\newcommand{\ceil}[1]{\lceil{#1}\rceil}
\newcommand{\floor}[1]{\lfloor{#1}\rfloor}
\newcommand{\rarr}[0]{\rightarrow}
\newcommand{\IM}[1]{\text{Im}(#1)}
\newcommand{\Hom}[0]{\text{Hom}}
\newcommand{\Pow}[1]{\text{Pow}(#1)}
\newcommand{\angles}[1]{\langle #1 \rangle}

%End custom commands
%Begin custom theorems

\newtheorem{innercustomgeneric}{\customgenericname}
\providecommand{\customgenericname}{}
\newcommand{\newcustomtheorem}[2]{%
\newenvironment{#1}[1]
{%
\renewcommand\customgenericname{#2}%
\renewcommand\theinnercustomgeneric{##1}%
\innercustomgeneric%
}
{\endinnercustomgeneric}
}

\newcustomtheorem{THRM}{Theorem}
\newcustomtheorem{LEM}{Lemma}
\newcustomtheorem{EX}{Exercise}
\newcustomtheorem{COR}{Corollary}
\newcustomtheorem{PROB}{Problem}
\newcustomtheorem{claim}{Claim}

\newenvironment{solution}
{\renewcommand\qedsymbol{$\blacksquare$}\begin{proof}[Solution]}
{\end{proof}}

\newenvironment{definition}
{\renewcommand\qedsymbol{$\blacksquare$}\begin{proof}[Definition]}
{\end{proof}}

%End custom theorems
\begin{document}

\title{Exercise 3}
\author{Arham Lodha}%put your name here
\date{March 28, 2025}%enter the due date here
\maketitle

\begin{PROB}{1a}\label{Problem 1a.}
\end{PROB}

\begin{PROB}{1b}\label{Problem 1b.}
\end{PROB}
\begin{proof}
    $\forall a, d, b, c \in \Z \ni ad - bc = 1$ and $\forall z \in \HC$. 

    Note the following, $\forall (m, n) \in \Z \setminus \left\{ \left(0,0\right) \right\}$ 

    \begin{align*}
        \left(\frac{1}{m\left(\frac{az + b}{cz + d}\right) + n}\right)^2 &= \left(\frac{1}{\frac{m(az + b) + n(cz + d)}{cz + d}}\right)^2 \\
        &= \left(\frac{cz + d}{z(ma + nc) + (mb + nd)}\right)^2
    \end{align*}

    Analogously

    \[\frac{1}{\abs{m\left(\frac{az + b}{cz + d}\right) + n}^{2\epsilon}} = \frac{\abs{cz + d}^{2 \epsilon}}{\abs{z(ma + nc) + (mb + nd)}^{2 \epsilon}}\]

    Thus,

    \begin{align*}
        G_{2, \epsilon}\left(\frac{az + b}{cz + d}\right) &= \frac{1}{2} \sum_{(m, n) \neq (0,0)} \left(\frac{1}{m\left(\frac{az + b}{cz + d}\right) + n}\right)^2 \frac{1}{\abs{m\left(\frac{az + b}{cz + d}\right) + n}^{2\epsilon}} \\
        &= \frac{1}{2}(cz + d)^2\abs{cz + d}^{2 \epsilon} \sum_{(m, n) \neq (0,0)} \frac{1}{\abs{z(ma + nc) + (mb + nd)}^{2 \epsilon}} \frac{1}{\left(z(ma + nc) + (mb + nd)\right)^2} \\
        &= \frac{1}{2}(cz + d)^2\abs{cz + d}^{2 \epsilon} \sum_{(u, v) \neq (0,0)} \frac{1}{\abs{uz + v}^{2 \epsilon}} \frac{1}{(u z + v)^2} \\
        &= (cz + d)^2\abs{cz + d}^{2 \epsilon} G_{2, \epsilon}(z)
    \end{align*}

    Where the 2nd last equality comes from the fact that $\begin{bmatrix}
        a & c \\ b & d
    \end{bmatrix} \in SL_2(\Z)$ is a invertible linear transformation which takes every nonzero lattice element of $\Z^2$ to a unique nonzero lattice element of $\Z^2$.  
\end{proof}

\bigskip

\begin{PROB}{2a}\label{Problem 2a.}
\end{PROB}
\begin{proof}
    \begin{enumerate}
        \item $Ix = x$:  \[I x = \frac{1x}{1} = x\]
        \item $\forall a_{i}, b_{i}, c_{i}, d_{i} \in \Z$ for $i \in [1, 2]$ where $a_id_i - b_ic_i = 0$. $\forall x \in P^1(\Q)$ 
        
        \begin{align*}
            \left(\begin{bmatrix}
                a_1 & b_1 \\ c_1 & d_1
            \end{bmatrix} \begin{bmatrix}
                a_2 & b_2 \\ c_2 & d_2
            \end{bmatrix}\right) (x) &= \begin{bmatrix}
                a_1a_2 + b_1c_2 & a_1b_2 + b_1 d_2 \\
                c_1a_2 + d_1c_2 & c_1b_2 + d_1d_2
            \end{bmatrix} (x) \\
            &= \frac{x(a_1a_2 + b_1c_2) + a_1b_2 + b_1d_2}{x(c_1a_2+d_1c_2) + c_1b_2+d_1d_2} \\
            &= \frac{a_1(a_2x + b_2) + b_1(c_2x + d_2)}{c_1(a_2x + b_2) + d_1(c_2x + d_2)} \\
            &= \frac{a_1\frac{a_2x + b_2}{c_2x + d_2} + b_1}{c_1\frac{a_2x + b_2}{c_2x + d_2} + d_1} \\
            &= \begin{bmatrix}
                a_1 & b_1 \\ c_1 & d_1
            \end{bmatrix}\left(\frac{a_2x + b_2}{c_2x + d_2}\right) \\
            &= \begin{bmatrix}
                a_1 & b_1 \\ c_1 & d_1
            \end{bmatrix}\left( \begin{bmatrix}
                a_2 & b_2 \\ c_2 & d_2
            \end{bmatrix} x\right)
        \end{align*}
    \end{enumerate}
    
    Thus the mobius transform is a valid action on $P^1(\Q)$ 
\end{proof}

\begin{PROB}{2b}\label{Problem 2b.}
\end{PROB}
\begin{proof}
    It suffices to show that $\forall x \in P^1(\Q)$, $\exists g \in SL_2(\Z) \ni g \cdot 0 = x$.  
    
    \begin{enumerate}
        \item $0 \to \infty$: Let $g = \begin{bmatrix}
            0 & -1 \\
            1 & 0
        \end{bmatrix} \implies g \cdot 0 = \frac{-1}{1*0} = \infty$
        \item $0 \to x$ for $x \in \Q \setminus \left\{ 0 \right\}$: $x = \frac{p}{q}$ for $p, q\in \Z \setminus \left\{ 0 \right\} \ni (p, q) = 1$. Thus $\exists m, n \in \Z \ni mp - nq = 1$. Let $g = \begin{bmatrix}
            -n & p \\ -m & q
        \end{bmatrix} \in SL_2(\Z) \implies g \cdot 0 = \frac{-n(0) + p}{-n(0) + q} = \frac{p}{q}$ 
        \item Note $0 \to 0$ is trivial since the mobius transform is an action.       
    \end{enumerate} 
    
    Thus the action of $SL_2(\Z)$ on $P^1(\Q)$ is transitive since everything in $P^1(\Q)$ is in the orbit of $0$. 
    
    
    \textbf{TODO}: Proof for finite orbits is trivial use finite index of $\Gamma$.  
\end{proof}

\end{document}